% Captions for the 14 November consolidation figures. Manuscript is unchanged.
% s is the number of completed optimizer updates. Requires amsmath/amssymb.
\newcommand{\consolidationGDCurvatureCaption}{%
Training loss and curvature diagnostics under full-batch GD. Columns identify
dataset/model; colors identify learning rate. Top: unsmoothed full-objective
loss on logarithmic axes. Bottom: $\eta\widehat\kappa_s$ (solid) and
$\eta\lambda_{\max}$ (dashed), both smoothed with the same centered 101-step
arithmetic mean; the dotted line marks two. Spectra are measured every 100
steps and linearly interpolated before smoothing; edge windows are truncated
equally. The horizontal axis counts completed optimizer updates. The gradient
diagnostic uses $g_{s-1},g_s$, and spectra retain their measurement-state index.
Spectral objectives use all 5,000 MLP images, 5,000/50,000 ResNet images and
2,500/59,680 Transformer sequences. One initialization per model. Smoothing is
retrospective display processing, not an estimate of exact transition timing.}

\newcommand{\consolidationGDGeometryCaption}{%
Consecutive gradient cosine, gradient norm and normalized displacement under
GD. Cosine and norm use complete trailing 50-observation means; norms use
logarithmic axes. $R_s$ is unsmoothed and includes all updates completed at step
$s$. Gradient statistics are indexed by the state where the gradient is
evaluated. Colors and models match the GD curvature figure. One initialization
per model; these diagnostics describe directional interaction and accumulation,
not generalization or a stability guarantee.}

\newcommand{\consolidationBackwardCaption}{%
Backward gradient alignment at completed step $s=2256$. Each point is
$\cos(g_s,g_{s-k})$ for a fixed endpoint. Top: lags 1--256; bottom: lags 1--32.
Learning rates are 0.04, 0.1 and 0.2 for MLP, ResNet and Transformer.
These separately recomputed GD continuations start at state 2000 and use one
initialization, without smoothing or endpoint averaging. The horizontal
coordinate is backward separation in steps, not training progress.}

\newcommand{\consolidationBatchCaption}{%
Mean sampled-gradient cosine over 199 adjacent pairs during 200-update
continuations, started at initialization or after 2,000 GD updates. Lines
average repetitions within initialization and weight three initialization
blocks equally; shading spans their range. Hollow rightmost markers denote
full GD. ResNet batches contain complete fixed 1,000-image normalization groups;
Transformer batches count prepared sequences. Batch axes are logarithmic.
Equal update counts do not imply equal data exposure.}

\newcommand{\consolidationAdamOverviewCaption}{%
High-rate full-batch Adam trajectories on CIFAR-10/MLP ($\eta=0.001$),
CIFAR-10/ResNet-32 (0.003) and WikiText-2/Transformer (0.003). Rows show training
loss, projected raw-gradient change $\eta\widehat\kappa_s$, consecutive
raw-gradient cosine and gradient norm relative to its initial value. Filled
gray bands mark pre-rise-baseline-to-peak intervals of recovered excursions
whose peak exceeds baseline loss by more than 10\%; outlined bands mark such
excursions with recovery unobserved. Loss and relative norm use logarithmic
axes. Projected-change axes start at zero, leaving 13/17/1 negative observations
outside the view; values are unmodified. Loss covers steps 0--1024, norms
0--1023 and paired diagnostics 1--1023. One initialization per model. Under
Adam, projected gradient change is not a validated Hessian estimate or
universal stability threshold.}

\newcommand{\consolidationTracingCaption}{%
Loss and functional motion across three unchanged-Adam continuations. Columns
show MLP, ResNet and Transformer; rows show full/probe loss, cosine between
successive representation changes and displacement coordinate RMS divided by
each stage's own window maximum. Heatmap rows are measured stage endpoints.
Each change $A_s-A_{s-1}$ is placed at completed step $s$, alongside the loss at
that state, without smoothing. Blue dashed lines mark the first observed
negative output-change cosine within the display; MLP is already negative at
entry, so its onset is unobserved. Black dashed lines mark the first full loss
at least 10\% above the minimum preceding the largest displayed peak. Both
markers continue through all rows. Fixed probes contain 5,000 images, 1,000
images and 2,240 token positions. MLP probe and full loss coincide; the other
probes are subsets. Displayed steps are 451--501, 390--450 and 553--601.
ResNet's selected excursion rises 311.53\% from step 404 to peak 415 and
recovers at 439. It exactly matches the native control used for the previous
trace, rather than bitwise replay of the original full-run overview.
Normalized RMS colors compare time within a stage, not absolute magnitudes
across stages. These are selected descriptive windows from one initialization
each. The temporal-structure companion now uses the same windows; its ResNet
SVD is freshly fitted to replay-qualified captures for steps 390--450.}

\newcommand{\consolidationTemporalStructureCaption}{%
Temporal structure of the same windows as the loss/motion figure: completed
steps 451--501 (MLP), 390--450 (ResNet) and 553--601 (Transformer), with
51/61/49 fitted displacement rows. The ResNet SVD is freshly fitted to the
replay-qualified larger excursion, peaking at 415 and recovering at 439.
Top: cumulative singular-value energy for output changes and raw gradients;
the dotted line marks 95\%. Middle: first-mode energy, first-ten-mode energy
and top-ten temporal subspace overlap with output at the selected stages.
Bottom: the leading output coefficient divided by the square root of its
coordinate count. Each coefficient uses one fixed pattern and sign convention
over its stated window; representation changes $A_s-A_{s-1}$ are placed at
completed step $s$. Raw-gradient rows are the undifferenced $g_{s-1}$ used by
the update into $s$. Predecessor rows for adjacent cosine are excluded from
the fits. Outputs are raw linear values for MLP and log probabilities for
ResNet and Transformer. No temporal centering, example averaging or row
normalization is applied to the primary SVD. Temporal overlap compares
time-domain subspaces, not identical neuron directions. Different probes and
window lengths preclude a controlled intrinsic-dimension comparison. Spectral
concentration alone is time-permutation invariant and establishes neither
causal origin nor spike specificity.}
